\documentclass[class=NCU_thesis, crop=false]{standalone}
\begin{document}

\chapter{緒論}

\section{研究動機}

根據衛生福利部統計資料 ~\cite{mohw2026report} ，台灣登記在案之聽覺機能障礙人口已超過十四萬人，
台灣手語（Taiwan Sign Language, TSL） \cite{tsl_website}為其日常溝通的重要媒介。
然而，目前多數手語學習與查詢資源仍以單詞對應影片之形式呈現，使用者在表達完整句子時，
需逐詞查找並自行組合，導致資訊呈現零散且缺乏連續性，難以有效支援自然語言溝通需求。
近年來，手語翻譯（Sign Language Translation, SLT）技術已能將自然語言轉換為手語詞彙序列（gloss sequence），
為文字到手語的轉換提供初步解決方案 \cite{Chu2025SLT} \cite{abdullah2025stateofthearttranslationtexttoglossusing} 。
然而，如何將多個手語詞彙進一步轉換為連續且自然的手語動作，仍是一項具挑戰性的問題。

在虛擬角色動作生成方面，傳統方法多依賴真人動作捕捉（motion capture）設備以取得高精度動作資料，
例如 Vicon、OptiTrack、Qualisys 等光學式動捕系統 \cite{optitrack_mocap} \cite{qualisys_mocap}，或 Xsens 等穿戴式慣性動捕系統 \cite{xsens_mocap}。
此類系統通常需搭配專用相機、感測器、軟體與校正流程，雖可取得高精度人體姿態資料，
但建置成本與操作門檻較高，不利於大規模擴充手語詞彙資料。
相較之下，直接利用既有手語影片資料庫作為動作來源，具備成本低與資料取得容易等優勢，
但亦面臨不同影片之間動作風格不一致、片段難以直接串接等問題，特別是在多單詞手語序列中，
容易產生動作不連續與不自然的現象。

此外，中文至 gloss 轉換、手語影片追蹤資料與虛擬人渲染模組分別使用不同的資料表示方式。
若缺乏明確的詞彙索引、資料交換格式與完整性檢查，即使各模組能獨立運作，仍難以形成可重複執行的完整動畫生成流程。
因此，如何將語言轉換結果對應至可查詢、可串接的手語動作資料，並使串接結果能直接供虛擬人渲染模組使用，亦是本研究欲處理的系統整合問題。

因此，本研究動機在於整合自然語言處理之手語翻譯模組與基於影片資料的動作生成方法，
建立一套由文字輸入至虛擬人手語輸出的模組化系統流程，並針對手語動作資料建置、多個手語動作片段之串接及模組間資料銜接提出設計方法，
以提升手語動畫之流暢性與自然度，進而改善手語數位內容之可用性與實用價值。


\pagebreak

\section{研究目的}

本研究旨在建構一套結合自然語言處理與虛擬角色動畫生成技術之手語生成系統，使使用者能夠透過自然語言輸入，自動產生具連續性且可觀看之虛擬人手語動畫，以提升手語數位應用之可用性與實用價值。具體研究目標如下：

\begin{itemize}
    \item 建立一套模組化之中文至虛擬人手語動畫生成流程，定義語言轉換、動作資料檢索、參數串接與虛擬人渲染之處理順序及資料交換方式。
    
    \item 整合既有之手語翻譯模組（Sign Language Translation, SLT） \cite{Chu2025SLT}、EHM-Tracker 與虛擬人動作生成系統（GUAVA） \cite{zhang2025guava}，分別處理文字至手語詞彙序列（gloss sequence）轉換、影片動作參數化及虛擬人動畫生成。
    
    \item 將既有手語詞彙影片建置為可由 gloss 查詢、可重複載入且可進行串接的參數化動作資料，並於動畫生成前檢查詞彙與追蹤資料之完整性。
    
    \item 針對多個手語詞彙動作片段串接時所產生之不連續問題，導入 motion envelope 有效區間擷取、詞彙層級參數平滑與過渡幀生成，以改善動作之間的連續性。

    \item 透過詞彙平滑度指標、EHM-Tracker 最佳化步數比較、系統整合展示與使用者問卷，評估系統生成結果之平滑度、追蹤品質、可觀看性及目前限制。
\end{itemize}


\pagebreak



\section{論文架構}

本論文共分為五個章節，內容安排如下：

\begin{itemize}
\item \textbf{第一章、緒論}：說明本研究之研究背景、研究動機與研究目標，並介紹論文整體架構。

\item \textbf{第二章、文獻回顧}：探討手語翻譯技術、虛擬角色動作生成方法以及人體動作建模相關研究，並分析現有方法於手語動作連續性方面之限制。

\item \textbf{第三章、研究方法}：詳細說明本研究之系統架構與系統設計，包括手語動作資料建置、使用者動畫生成流程、模組介面與例外處理，以及手語翻譯、動作追蹤、片段串接與虛擬角色生成方法。

\item \textbf{第四章、實驗設計與結果}：呈現系統介面與整合輸出，並針對 motion envelope 區間判定、詞彙層級參數平滑、EHM-Tracker 最佳化步數及使用者評估結果進行分析與討論。

\item \textbf{第五章、結論與未來展望}：總結本研究之主要成果與貢獻，說明系統限制，並探討手語翻譯與動作生成過程中仍待改進之問題與未來研究方向。
\end{itemize}

\pagebreak

\end{document}
